\documentclass[10pt]{article}
\usepackage[margin=2cm]{geometry}
\usepackage{amsmath,amssymb,amsfonts,booktabs,array,multirow}
\usepackage[table]{xcolor}
\renewcommand{\thetable}{S\arabic{table}}
\begin{document}
\begin{center}{\Large Supplementary Material}\\[4pt]
Physics-Guided Probabilistic Glucose Forecasting in Type 1 Diabetes: An Audited Multi-Horizon Evaluation with Validated Patient-Specific Insulin Sensitivity\end{center}

\section*{Notation, Variables, and Terminology}
Rate constants are expressed per minute unless stated otherwise. A subscript $b$ denotes a basal or fasting quantity, a hat denotes a prediction, bold lowercase symbols denote vectors, and bold uppercase symbols denote matrices. Bracketed reference numbers refer to the reference list of the main text. The patient-parameter vector $\theta_p$ is listed in Table~XIV of the main text.

\begin{table}[h!]
\centering
\caption{Physiological state variables and external inputs}
\label{tab:states}
\footnotesize
\begin{tabular}{clp{5cm}p{7cm}}
\toprule
Symbol & Unit & Name & Meaning in this study \\
\midrule
$G(t)$ & mg/dL & Plasma glucose & Glucose concentration forecast by the learned spline and constrained by the mechanistic equation. \\
$G_b$ & mg/dL & Basal glucose & Subject-specific fasting anchor used by the Bergman equilibrium. \\
$G_0$ & mg/dL & Anchor glucose & Last real CGM observation at forecast time $t=0$. \\
$X(t)$ & min$^{-1}$ & Remote insulin action & Interstitial insulin effect entering glucose disposal additively with $p_1$. \\
$I(t)$ & $\mu$U/mL & Plasma insulin & Circulating insulin concentration generated by subcutaneous absorption and clearance. \\
$I_b$ & $\mu$U/mL & Basal plasma insulin & Concentration implied by the subject's basal delivery; subtracting it preserves basal equilibrium. \\
$S_1, S_2$ & U & SC insulin states & First and second subcutaneous absorption compartments. \\
$Q_{\text{sto}}$ & g & Stomach carbohydrate & Ingested carbohydrate awaiting gastric emptying. \\
$Q_{\text{gut}}$ & g & Gut carbohydrate & Carbohydrate available for intestinal absorption. \\
$R_a(t)$ & mg/min & Glucose appearance rate & Meal-derived glucose entering the accessible glucose compartment. \\
$u_{\text{ins}}$ & U/min & Insulin input & Basal plus bolus insulin delivery on the five-minute grid. \\
$u_{\text{carb}}$ & g/min & Meal input & Recorded carbohydrate ingestion on the grid. \\
$\mathbf{z}$ & -- & Compartment state & ${[S_1, S_2, I, Q_{\text{sto}}, Q_{\text{gut}}, X]}^{\mathsf T}$. \\
$\mathbf{u}$ & -- & Input vector & ${[u_{\text{ins}}, u_{\text{carb}}, I_b]}^{\mathsf T}$. \\
\bottomrule
\end{tabular}
\end{table}

\begin{table}[h!]
\centering
\caption{Architecture, discretization, and uncertainty notation}
\label{tab:arch}
\footnotesize
\begin{tabular}{lll}
\toprule
Symbol & Value or dimension & Definition \\
\midrule
$t$ & 0--120 min & Time measured from the forecast anchor. \\
$h$ & 30, 60, 90, 120 min & Forecast horizon. \\
$\Delta$ & 5 min & Common temporal grid step. \\
$B$ & batch dependent & Number of windows processed together. \\
$H$ & 4 & Number of reported forecast horizons. \\
$K$ & 12 & Number of cubic B-spline basis functions. \\
$B_k(t)$ & spline basis & The $k$th clamped, uniformly knotted B-spline basis function. \\
$c_k$ & learned scalar & Spline coefficient emitted by the trajectory head. \\
$G_\theta(t)$ & mg/dL & Anchored learned trajectory $G_0 + \sum_k c_k[B_k(t) - B_k(0)]$. \\
$A, B$ & $6\times6$, $6\times3$ & Continuous compartment-state and input matrices. \\
$A_d, B_d$ & discrete matrices & Exact zero-order-hold matrices obtained from one augmented matrix exponential. \\
$\theta_p$ & 10 elements & Patient-specific physiological parameter vector in Table~XIV of the main text. \\
$g = \sigma(\gamma)$ & initialized near 0.10 & Learned trust gate controlling the mechanistic prior. \\
$\hat{G}_{0.5}$ & mg/dL & Median forecast and reported point prediction. \\
$\hat{G}_q$ & mg/dL & Forecast at quantile level $q \in \{0.1, 0.5, 0.9\}$. \\
$\delta_q(h)$ & mg/dL & Positive softplus-derived distance from the median, preventing quantile crossing. \\
$r(t)$ & dimensionless & Scale-normalized physics residual evaluated at 121 collocation points. \\
$n_{\text{eff}}$ & $\approx N/48$ & Approximate independent sample count after accounting for window overlap. \\
\bottomrule
\end{tabular}
\end{table}

\begin{table}[h!]
\centering
\caption{Training-loss terminology}
\label{tab:loss}
\footnotesize
\begin{tabular}{lp{13cm}}
\toprule
Term & Meaning \\
\midrule
$\mathcal{L}_{\text{Huber}}$ & Horizon-weighted robust point loss for the median forecast; quadratic near zero error and linear for large errors. \\
$\mathcal{L}_{\text{pinball}}$ & Proper quantile loss $\max\{q(y-\hat{y}_q), (q-1)(y-\hat{y}_q)\}$. \\
$\mathcal{L}_{\text{res}}$ & Mean squared, nondimensionalized Bergman residual over the collocation grid. \\
$\mathcal{L}_{\text{prior}}$ & Weak population-prior penalty that discourages physiologically implausible parameter drift. \\
$\mathcal{L}_{\text{temporal}}$ & Within-subject penalty on disagreement between parameter estimates from adjacent windows. \\
$w_{\text{phys}}$ & Physics weight, adaptive in the relevant ablation arms and optionally multiplied by a curriculum ramp. \\
$s_i = \log\sigma_i^2$ & Stable log-variance parameterization used for learned multi-objective weighting~[51]. \\
\bottomrule
\end{tabular}
\end{table}

\begin{table}[h!]
\centering
\caption{Evaluation metrics and clinical terms}
\label{tab:metrics}
\footnotesize
\begin{tabular}{llp{12.3cm}}
\toprule
Term & Unit & Definition and interpretation \\
\midrule
MAE & mg/dL & Mean absolute error, $n^{-1}\sum_i |\hat{y}_i - y_i|$. \\
RMSE & mg/dL & Root-mean-square error, emphasizing larger deviations. \\
$R^2$ & -- & $1 - \text{SS}_{\text{res}}/\text{SS}_{\text{tot}}$. \\
MARD & \% & Mean absolute relative difference, with reference glucose in the denominator. \\
Skill & \% & $1 - \text{RMSE}_{\text{model}}/\text{RMSE}_{\text{persistence}}$. Positive values beat no-change prediction. \\
TIR/TAR/TBR & \% & Time in, above, or below the consensus glucose range. The main in-range interval is 70--180 mg/dL~[55]. \\
CV ratio & -- & $\text{CV}_{\text{pred}}/\text{CV}_{\text{actual}}$; values below one quantify excursion compression. \\
LBGI/HBGI & -- & Low/high blood-glucose risk indices after the Kovatchev symmetric risk transform. \\
Coverage & -- & Fraction of observations below a stated predictive quantile; calibration requires observed and nominal coverage to agree. \\
Hypo sensitivity & -- & Fraction of actual $G<70$ mg/dL events detected by the forecast or alarm. \\
Precision & -- & Fraction of hypoglycemia alarms corresponding to actual events. \\
ICC(1) & -- & One-way, single-measure intraclass correlation used for $S_I$ test--retest stability~[60]. \\
$\rho$ & -- & Spearman rank correlation used for monotonic association with insulin dose~[61]. \\
$\text{IG}_i$ & mg/dL & Integrated-gradient attribution of feature $i$ relative to the standardized-mean baseline. \\
Clarke/Parkes EGA & zone & Point-wise clinical error grids originally designed for glucose-measurement error. \\
PRED-EGA & category & Prediction-specific error grid combining level and rate information~[69]; reported in the Results of the main text. \\
\bottomrule
\end{tabular}
\end{table}

\begin{table}[h!]
\centering
\caption{Abbreviations and protocol terminology}
\label{tab:abbrev}
\footnotesize
\begin{tabular}{lp{4.5cm}p{9cm}}
\toprule
Short form & Expansion & Study-specific meaning \\
\midrule
T1D/IDDM & Type 1 diabetes / insulin-dependent diabetes mellitus & IDDM is retained only when describing older physiological literature. \\
CGM & Continuous glucose monitoring & Nominal five-minute glucose sampling in OhioT1DM. \\
SC & Subcutaneous & Route represented by insulin compartments $S_1$ and $S_2$. \\
IOB/COB & Insulin/carbohydrate on board & Derived as $S_1+S_2$ and $Q_{\text{sto}}+Q_{\text{gut}}$. \\
TDD/ICR & Total daily insulin dose / insulin-to-carbohydrate ratio & External consistency variables for the physiology analysis; ICR is available only for six subjects. \\
PINN/ODE & Physics-informed neural network / ordinary differential equation & Neural training constrained by the physiological differential equations. \\
LOSO & Leave-one-subject-out & Twelve subject-disjoint folds, each withholding all data from one test subject. \\
Official protocol & OhioT1DM temporal holdout & Later data from subjects whose earlier data are available for training; therefore personalized. \\
BGLP & Blood Glucose Level Prediction challenge & OhioT1DM challenge setting; 2018 and 2020 have different test warmup rules. \\
EGA & Error-grid analysis & Clarke and Parkes clinical zone classification. \\
IG/SHAP & Integrated gradients / SHapley Additive exPlanations & IG is used; the legacy single-timestep SHAP procedure was rejected. \\
SD/CI/ICC & Standard deviation / confidence interval / intraclass correlation & Uncertainty and reliability summaries computed at the subject level. \\
A0--A7 & Ablation arms & Controlled variants of physics form, weighting, parameter personalization, and curriculum. \\
ZOH & Zero-order hold & Exact interval propagation under piecewise-constant inputs; also used in the literature for persistence. \\
\bottomrule
\end{tabular}
\end{table}
\end{document}
